\documentclass[12pt]{article}

% === Кодировки и язык ===
\usepackage{cmap}
\usepackage[T2A]{fontenc}
\usepackage{mathtext}
\usepackage[utf8]{inputenc}
\usepackage[english, russian]{babel}

% === Цвета ===
\usepackage[dvipsnames]{xcolor}

% === Математика ===
\usepackage{amsmath, amssymb, wasysym, amsthm}
\usepackage{stmaryrd}
\usepackage{proof}

% === Графика и float'ы ===
\usepackage{wrapfig, float}
\usepackage{graphicx}

% === TikZ / PGFPlots ===
\usepackage{pgfplots}
\pgfplotsset{compat=newest}
\usepackage{pgfplotstable}
\usepackage{tikz, tikz-3dplot}
\usetikzlibrary{calc, patterns, angles, quotes, intersections, automata, positioning, arrows.meta}

% === Таблицы ===
\usepackage{tabularx}
\usepackage{xltabular}
\renewcommand\tabularxcolumn[1]{m{#1}}

% === Разметка ===
\usepackage{geometry}
\geometry{margin=0.5in}

% === Списки и прочее ===
\usepackage[shortlabels]{enumitem}
\usepackage{verbatim}

% ==========================================
% === ТЕОРЕМНЫЕ ОКРУЖЕНИЯ (amsthm) ===
% ==========================================
\theoremstyle{plain}
\newtheorem{theorem}{Теорема}[section]
\newtheorem{corollary}{Следствие}[section]
\newtheorem{lemma}{Лемма}[section]
\newtheorem{statement}{Утверждение}[section]
\newtheorem{axiom}{Аксиома}[section]
\newtheorem{law}{Закон}[section]

\theoremstyle{definition}
\newtheorem{definition}{Определение}[section]
\newtheorem{example}{Пример}[section]
\newtheorem{problem}{Задача}
\newtheorem{method}{Метод}[section]
\newtheorem{event}{Событие}[section]

\theoremstyle{remark}
\newtheorem{note}{Заметка}[section]
\newtheorem{property}{Свойство}[section]
\newtheorem{properties}{Свойства}[section]
\newtheorem{principle}{Правило}[section]
\newtheorem{person}{Личность}[section]
\newtheorem*{solution}{Решение}

% ==========================================
% === КРАСИВЫЕ РАМКИ (tcolorbox) ===
% ==========================================
\usepackage{tcolorbox}
\tcbuselibrary{skins, breakable, theorems}

% Базовые стили для разных типов блоков
\tcbset{
    % Стиль для теорем, лемм (Синий - строгий)
    thmstyle/.style={
        enhanced, breakable,
        frame hidden, boxrule=0pt,
        colback=blue!5!white, % Фон: 5% синего, 95% белого
        borderline west={2.5pt}{0pt}{blue!70!black}, % Вертикальная линия слева
        before skip=10pt, after skip=10pt,
        left=4pt, right=4pt, top=4pt, bottom=4pt
    },
    % Стиль для доказательств (фиолетовый - процесс рассуждения)
    proofstyle/.style={
        enhanced, breakable,
        frame hidden, boxrule=0pt,
        colback=Plum!5!white,
        borderline west={2.5pt}{0pt}{Plum!80!black},
        before skip=10pt, after skip=10pt,
        left=4pt, right=4pt, top=4pt, bottom=4pt,
        sharp corners
    },
    % Стиль для определений и аксиом (Зеленый - фундаментальный)
    defstyle/.style={
        enhanced, breakable,
        frame hidden, boxrule=0pt,
        colback=ForestGreen!5!white,
        borderline west={2.5pt}{0pt}{ForestGreen!80!black},
        before skip=10pt, after skip=10pt,
        left=4pt, right=4pt, top=4pt, bottom=4pt
    },
    % Стиль для примеров и задач (Оранжевый - практический)
    exstyle/.style={
        enhanced, breakable,
        frame hidden, boxrule=0pt,
        colback=orange!5!white,
        borderline west={2.5pt}{0pt}{orange!90!black},
        before skip=10pt, after skip=10pt,
        left=4pt, right=4pt, top=4pt, bottom=4pt
    },
    % Стиль для заметок и свойств (Серый - нейтральный)
    remstyle/.style={
        enhanced, breakable,
        frame hidden, boxrule=0pt,
        colback=gray!5!white,
        borderline west={2.5pt}{0pt}{gray!70!black},
        before skip=10pt, after skip=10pt,
        left=4pt, right=4pt, top=4pt, bottom=4pt
    }
}

% Применяем стили к существующим окружениям amsthm
% 1. Строгая математика
\tcolorboxenvironment{theorem}{thmstyle}
\tcolorboxenvironment{lemma}{thmstyle}
\tcolorboxenvironment{corollary}{thmstyle}
\tcolorboxenvironment{statement}{thmstyle}
\tcolorboxenvironment{law}{thmstyle}

% 2. Доказательство
\tcolorboxenvironment{proof}{proofstyle}
\renewcommand{\qed}{}

% 3. Определения
\tcolorboxenvironment{definition}{defstyle}
\tcolorboxenvironment{axiom}{defstyle}

% 4. Практика
\tcolorboxenvironment{example}{exstyle}
\tcolorboxenvironment{problem}{exstyle}
\tcolorboxenvironment{method}{exstyle}
\tcolorboxenvironment{event}{exstyle}

% 5. Заметки и прочее
\tcolorboxenvironment{note}{remstyle}
\tcolorboxenvironment{property}{remstyle}
\tcolorboxenvironment{properties}{remstyle}
\tcolorboxenvironment{principle}{remstyle}
\tcolorboxenvironment{person}{remstyle}
\tcolorboxenvironment{solution}{remstyle}

% === Гиперссылки (почти последним!) ===
\usepackage[
colorlinks=true,
linkcolor=blue,
urlcolor=blue,
citecolor=blue,
bookmarks=true,
bookmarksopen=false
]{hyperref}


\begin{document}
	\tableofcontents
	\setcounter{tocdepth}{4}
	\newpage
	
	\begin{center}
		\textbf{\Large Ответы на вопросы к зачёту по обществознанию}
	\end{center}
	
	\section{Гражданское общество и правовое государство}
	
	\subsection{Гражданское общество}
	
	\textbf{Гражданское общество}~--- это сфера общественной жизни, которая не контролируется государством, где люди действуют независимо и самостоятельно.
	
	\textbf{Элементы гражданского общества:}
	\begin{itemize}
		\item Независимые СМИ~--- обеспечивают свободу информации и контроль над властью.
		\item Частный бизнес~--- гарантирует экономическую независимость граждан.
		\item Органы местного самоуправления~--- помогают гражданам осуществлять власть на местах.
		\item Общественные организации (НКО, профсоюзы, правозащитные организации).
		\item Политические партии и движения.
	\end{itemize}
	
	\textbf{Пример:} Общество защиты прав потребителей~--- организация, готовая вступиться за тех, чьи права были нарушены.
	
	\textbf{Значение:} Гражданское общество является основой демократии, так как обеспечивает участие граждан в управлении государством и защиту их интересов.
	
	\subsection{Правовое государство}
	
	\textbf{Правовое государство}~--- это государство, в котором исключена возможность злоупотребления властью, а все граждане и органы власти равны перед законом.
	
	\textbf{Принципы правового государства:}
	\begin{enumerate}
		\item \textbf{Верховенство права и закона}~--- закон обязателен для всех, включая государственные органы.
		\item \textbf{Незыблемость прав и свобод человека}~--- права человека признаются высшей ценностью.
		\item \textbf{Равенство всех перед законом}~--- отсутствие привилегий для отдельных лиц или групп.
		\item \textbf{Взаимная ответственность человека и государства}~--- государство несёт ответственность перед гражданами, а граждане~--- перед государством.
		\item \textbf{Разделение властей}~--- законодательная, исполнительная и судебная власти независимы друг от друга.
	\end{enumerate}
	
	\textbf{Вывод:} Демократия невозможна без правового государства, а правовое государство невозможно без развитого гражданского общества.
	
	\section{СМИ в политической системе общества}
	
	\subsection{Понятие СМИ}
	
	\textbf{СМИ (средства массовой информации)}~--- средства тиражирования и распространения информации для широкой аудитории.
	
	\subsection{История развития СМИ}
	\begin{itemize}
		\item \textbf{XVIII век}~--- появление первых газет (печатные СМИ).
		\item \textbf{XX век}~--- развитие аудиовизуальных СМИ (радио и телевидение).
		\item \textbf{XXI век}~--- появление интернет-СМИ.
		\item Кинематограф и реклама также являются формами СМИ.
	\end{itemize}
	
	\subsection{Признаки СМИ}
	\begin{enumerate}
		\item Периодичность издания.
		\item Наличие постоянного названия.
		\item Массовость и общедоступность.
	\end{enumerate}
	
	\subsection{Политические функции СМИ}
	\begin{enumerate}
		\item \textbf{Информационная}~--- донесение политической информации до граждан.
		\item \textbf{Политическая социализация}~--- формирование политических взглядов и ценностей.
		\item \textbf{Критика и контроль}~--- контроль над деятельностью властей (<<четвёртая власть>>).
		\item \textbf{Мобилизационная}~--- побуждение граждан к политическому участию.
	\end{enumerate}
	
	\subsection{Положительные и отрицательные стороны СМИ}
	
	\textbf{Положительные:}
	\begin{enumerate}
		\item Информирование граждан о происходящих событиях.
		\item Раскрытие негативных явлений в обществе (коррупция, нарушения прав).
		\item Способствование развитию демократии.
	\end{enumerate}
	
	\noindent \textbf{Отрицательные:}
	\begin{enumerate}
		\item Коммерциализированность~--- формирование потребительского отношения.
		\item Не приучают думать~--- информация усваивается в готовом виде.
		\item Возможность манипулирования общественным мнением.
		\item Распространение недостоверной информации (фейки).
	\end{enumerate}
	
	\section{Федеративное устройство РФ}
	
	\subsection{Понятие федерации}
	
	\textbf{Федерация}~--- форма государственного устройства, при которой части государства (субъекты) обладают определённой политической самостоятельностью.
	
	Россия согласно Конституции провозглашается федеративным государством.
	
	\subsection{Субъекты Российской Федерации}
	
	В состав РФ входит \textbf{89 субъектов} (с 2022 года), которые делятся на следующие виды:
	\begin{itemize}
		\item \textbf{Республики} (22)~--- имеют свою конституцию и государственный язык.
		\item \textbf{Края} (9).
		\item \textbf{Области} (48).
		\item \textbf{Города федерального значения} (3): Москва, Санкт-Петербург, Севастополь.
		\item \textbf{Автономная область} (1): Еврейская АО.
		\item \textbf{Автономные округа} (4).
	\end{itemize}
	
	\subsection{Принципы федеративного устройства РФ}
	\begin{enumerate}
		\item \textbf{Государственная целостность}~--- территория РФ едина и неделима.
		\item \textbf{Единство системы государственной власти}.
		\item \textbf{Разграничение предметов ведения} между федеральным центром и субъектами.
		\item \textbf{Равноправие субъектов} во взаимоотношениях с федеральным центром.
		\item \textbf{Равноправие и самоопределение народов} в РФ.
	\end{enumerate}
	
	\subsection{Распределение полномочий}
	
	\textbf{Исключительное ведение Федерации:}
	\begin{itemize}
		\item Принятие и изменение Конституции.
		\item Оборона и безопасность.
		\item Внешняя политика.
		\item Денежная эмиссия.
		\item Федеральные налоги.
	\end{itemize}
	
	\noindent \textbf{Совместное ведение Федерации и субъектов:}
	\begin{itemize}
		\item Защита прав и свобод человека.
		\item Вопросы образования, здравоохранения.
		\item Природопользование и охрана окружающей среды.
	\end{itemize}
	
	\noindent \textbf{Ведение субъектов:}~--- всё, что не отнесено к ведению Федерации и совместному ведению.
	
	\section{Органы государственной власти и устройство РФ}
	
	\subsection{Понятие органов государственной власти}
	
	\textbf{Органы государственной власти}~--- организации, учреждения и должностные лица, предназначенные для управления государством в разных сферах.
	
	\noindent \textbf{Отличительные признаки:}
	\begin{enumerate}
		\item Образуются в законодательном порядке.
		\item Наделены определёнными полномочиями.
		\item Финансируются из государственного бюджета.
	\end{enumerate}
	
	\subsection{Классификация органов власти}
	
	\textbf{По типу власти:}
	\begin{itemize}
		\item \textbf{Законодательные}~--- Федеральное Собрание (Государственная Дума, Совет Федерации).
		\item \textbf{Исполнительные}~--- Правительство РФ.
		\item \textbf{Судебные}~--- Конституционный суд, Верховный суд.
	\end{itemize}
	
	\noindent \textbf{По уровню:}
	\begin{itemize}
		\item Высшие (федеральные).
		\item Региональные (субъектов РФ).
		\item Местные (муниципальные).
	\end{itemize}
	
	\subsection{Президент РФ}
	
	\textbf{Статус:} Глава государства, гарант Конституции, Верховный главнокомандующий ВС РФ.
	
	\noindent \textbf{Требования к кандидату:}
	\begin{itemize}
		\item Возраст не менее 35 лет.
		\item Проживание на территории РФ не менее 25 лет (с 2020 года).
		\item Отсутствие иностранного гражданства.
	\end{itemize}
	
	\noindent \textbf{Полномочия:}
	\begin{itemize}
		\item Право вето на законы и право законодательной инициативы.
		\item Назначение Председателя Правительства.
		\item Внесение кандидатур судей высших судов.
		\item Представление РФ в международных отношениях.
		\item Подписание и обнародование федеральных законов.
		\item Право помилования.
	\end{itemize}
	
	\noindent \textbf{Импичмент (ст. 93 Конституции):}
	\begin{enumerate}
		\item Обвинение выдвигает Государственная Дума.
		\item Решение принимает Совет Федерации (не менее 2/3 голосов).
	\end{enumerate}
	
	\subsection{Федеральное Собрание}
	
	\textbf{Государственная Дума (нижняя палата):}
	\begin{itemize}
		\item 450 депутатов, избираемых на 5 лет.
		\item Принимает федеральные законы.
		\item Даёт согласие на назначение Председателя Правительства.
		\item Объявляет амнистию.
	\end{itemize}
	
	\noindent \textbf{Совет Федерации (верхняя палата):}
	\begin{itemize}
		\item По 2 представителя от каждого субъекта + 30 от Президента = 208 человек.
		\item Одобряет федеральные законы.
		\item Утверждает изменение границ.
		\item Назначает судей высших судов.
	\end{itemize}
	
	\noindent \textbf{Процедура принятия закона:}
	\begin{enumerate}
		\item Рассмотрение в Думе (3 чтения).
		\item Принятие большинством голосов.
		\item Рассмотрение Советом Федерации.
		\item Подписание Президентом.
	\end{enumerate}
	
	\subsection{Разделение властей}
	
	\textbf{Горизонтальное}~--- распределение между ветвями власти (законодательная, исполнительная, судебная) на одном уровне.
	
	\textbf{Вертикальное}~--- распределение между уровнями власти (федеральный, региональный, местный).
	
	\section{Понятие экономики. Этапы становления экономической науки}
	
	\subsection{Понятие экономики}
	
	\textbf{Экономика} (от греч. <<искусство ведения домашнего хозяйства>>) употребляется в двух значениях:
	\begin{itemize}
		\item Как \textbf{хозяйственная деятельность} общества.
		\item Как \textbf{наука}, изучающая эту деятельность.
	\end{itemize}
	
	\subsection{Функции экономики как науки}
	\begin{enumerate}
		\item \textbf{Познавательная}~--- изучение экономических процессов и явлений.
		\item \textbf{Практическая}~--- применение знаний для решения хозяйственных задач.
		\item \textbf{Прогностическая}~--- определение перспектив развития экономики.
		\item \textbf{Методологическая}~--- разработка методов экономического анализа.
	\end{enumerate}
	
	\subsection{Этапы становления экономической науки}
	
	\textbf{1. Экономическая мысль Древнего мира:}
	\begin{itemize}
		\item \textbf{Аристотель}~--- основатель экономической науки, выделил понятия собственности и богатства, разграничил экономику и хрематистику.
	\end{itemize}
	
	\noindent \textbf{2. Средневековье:}
	\begin{itemize}
		\item Экономика не выделялась как самостоятельная наука.
		\item Господство натурального хозяйства.
	\end{itemize}
	
	\noindent \textbf{3. Новое время (XVI--XVII вв.):}
	\begin{itemize}
		\item \textbf{Меркантилизм}~--- богатство страны определяется количеством денег (золота, серебра), главный источник~--- внешняя торговля.
		\item \textbf{Физиократы}~--- богатство создаётся в сфере производства, прежде всего в сельском хозяйстве; природа~--- <<мать богатств>>.
	\end{itemize}
	
	\noindent \textbf{4. Классическая политэкономия (XVIII--XIX вв.):}
	\begin{itemize}
		\item \textbf{Адам Смит}~--- труд как источник богатства, <<невидимая рука рынка>>.
		\item \textbf{Давид Рикардо}~--- теория стоимости и распределения.
	\end{itemize}
	
	\noindent \textbf{5. Современная экономическая наука (XX--XXI вв.):}
	\begin{itemize}
		\item \textbf{Кейнсианство}~--- необходимость государственного регулирования экономики.
		\item \textbf{Монетаризм}~--- роль денежной политики.
		\item \textbf{Институционализм}~--- влияние институтов на экономику.
	\end{itemize}
	
	\section{Типы экономических систем. Их преимущества и недостатки}
	
	\textbf{Экономическая система}~--- способ организации хозяйственной жизни общества, определяющий ответы на основные вопросы экономики: \textit{что, как и для кого производить}.
	
	\subsection{Традиционная экономическая система}
	
	\textbf{Характеристики:}
	\begin{itemize}
		\item Низкий уровень производства, преобладание сельского хозяйства.
		\item Что и как производить определяется обычаями и традициями.
		\item Господство натурального хозяйства.
		\item Медленное развитие.
	\end{itemize}
	
	\noindent \textbf{Преимущества:}
	\begin{enumerate}
		\item Стабильность и предсказуемость.
		\item Высокое качество ремесленной продукции.
		\item Отсутствие экономических кризисов.
	\end{enumerate}
	
	\noindent \textbf{Недостатки:}
	\begin{enumerate}
		\item Отсутствие технического прогресса.
		\item Ограниченность производимых продуктов.
		\item Низкий уровень жизни большинства населения.
	\end{enumerate}
	
	\subsection{Рыночная экономическая система}
	
	\textbf{Характеристики:}
	\begin{itemize}
		\item Частная собственность на средства производства.
		\item Экономическая свобода производителей и потребителей.
		\item Конкуренция как движущая сила.
		\item Цены определяются рынком (спросом и предложением).
	\end{itemize}
	
	\noindent \textbf{Преимущества:}
	\begin{enumerate}
		\item Эффективное распределение ресурсов.
		\item Гибкая реакция на спрос.
		\item Стимулирование научно-технического прогресса.
		\item Разнообразие товаров и услуг.
	\end{enumerate}
	
	\noindent \textbf{Недостатки:}
	\begin{enumerate}
		\item Нестабильность (циклические кризисы).
		\item Социальное неравенство.
		\item Недостаточное производство общественных благ.
		\item Безработица.
	\end{enumerate}
	
	\subsection{Командная (плановая) экономическая система}
	
	\textbf{Характеристики:}
	\begin{itemize}
		\item Государственная собственность на средства производства.
		\item Централизованное планирование.
		\item Государство определяет что, как и для кого производить.
		\item Фиксированные цены.
	\end{itemize}
	
	\noindent \textbf{Преимущества:}
	\begin{enumerate}
		\item Возможность концентрации ресурсов на важных направлениях.
		\item Отсутствие безработицы.
		\item Гарантированный минимальный уровень жизни.
		\item Отсутствие кризисов перепроизводства.
	\end{enumerate}
	
	\noindent \textbf{Недостатки:}
	\begin{enumerate}
		\item Отсутствие стимулов к качественному труду и инновациям.
		\item Дефицит товаров.
		\item Неэффективное распределение ресурсов.
		\item Отсутствие экономической свободы.
	\end{enumerate}
	
	\subsection{Смешанная экономическая система}
	
	Сочетает элементы рыночной и командной экономик. Рынок регулируется государством в целях социальной защиты и стабильности.
	
	\textbf{Национальные модели:}
	\begin{itemize}
		\item \textbf{Американская}~--- акцент на индивидуальную свободу и минимальное вмешательство государства.
		\item \textbf{Скандинавская (шведская)}~--- социально-ориентированная экономика с высокими налогами и развитой системой социальной защиты.
		\item \textbf{Японская (азиатская)}~--- государство как координатор, сочетание традиций и современности.
	\end{itemize}
	
	\section{Измерители экономической активности. Макроэкономические показатели}
	
	\subsection{Понятие макроэкономики}
	
	\textbf{Макроэкономика}~--- раздел экономической науки, изучающий экономику как целостную систему на уровне страны или мира.
	
	\subsection{Основные макроэкономические показатели}
	
	\textbf{1. ВВП (Валовой внутренний продукт):}
	\begin{itemize}
		\item Рыночная стоимость всех конечных товаров и услуг, произведённых на территории страны за определённый период (обычно год).
		\item Включает продукцию всех предприятий, находящихся на территории страны, независимо от их национальной принадлежности.
	\end{itemize}
	
	\noindent \textbf{2. ВНП (Валовой национальный продукт):}
	\begin{itemize}
		\item Рыночная стоимость всех конечных товаров и услуг, произведённых национальными предприятиями, независимо от их местоположения.
		\item ВНП = ВВП + доходы граждан за рубежом $-$ доходы иностранцев в стране.
	\end{itemize}
	
	\noindent \textbf{3. ЧВП (Чистый внутренний продукт):}
	\begin{itemize}
		\item ЧВП = ВВП $-$ амортизация (износ основных фондов).
	\end{itemize}
	
	\noindent \textbf{4. ЧНП (Чистый национальный продукт):}
	\begin{itemize}
		\item ЧНП = ВНП $-$ амортизация.
	\end{itemize}
	
	\noindent \textbf{5. Национальный доход:}
	\begin{itemize}
		\item Сумма всех доходов, полученных в экономике (зарплаты, прибыль, рента, процент).
	\end{itemize}
	
	\noindent \textbf{6. ВВП на душу населения:}
	\begin{itemize}
		\item Показатель уровня экономического развития страны.
		\item ВВП на душу = ВВП / численность населения.
	\end{itemize}
	
	\noindent \textbf{7. Другие показатели:}
	\begin{itemize}
		\item Уровень инфляции.
		\item Уровень безработицы.
		\item Индекс потребительских цен.
		\item Торговый баланс.
	\end{itemize}
	
	\section{Экономический цикл. Фазы экономического цикла}
	
	\subsection{Понятие экономического цикла}
	
	\textbf{Экономический цикл}~--- периодические колебания уровня экономической активности, чередование подъёмов и спадов.
	
	\subsection{Фазы экономического цикла}
	
	\textbf{1. Кризис (рецессия, спад):}
	\begin{itemize}
		\item Падение производства и ВВП.
		\item Рост безработицы.
		\item Снижение цен и доходов.
		\item Банкротство предприятий.
		\item Причина: чаще всего кризис перепроизводства.
	\end{itemize}
	
	\noindent \textbf{2. Депрессия (стагнация):}
	\begin{itemize}
		\item Прекращение падения, но отсутствие роста.
		\item Экономические показатели стабилизируются на низком уровне.
		\item Начинается обновление производства, внедрение новых технологий.
	\end{itemize}
	
	\noindent \textbf{3. Оживление (восстановление):}
	\begin{itemize}
		\item Постепенный рост производства.
		\item Снижение безработицы.
		\item Увеличение инвестиций.
		\item Достижение докризисного уровня.
	\end{itemize}
	
	\noindent \textbf{4. Подъём (экспансия, бум):}
	\begin{itemize}
		\item Активный рост всех экономических показателей.
		\item Рост занятости, доходов, инвестиций.
		\item Высшая точка~--- <<экономический бум>> или <<процветание>>.
		\item После достижения пика начинается новый спад.
	\end{itemize}
	
	\subsection{Исторические примеры кризисов}
	\begin{itemize}
		\item \textbf{Первый промышленный кризис}~--- Англия, XVIII век.
		\item \textbf{Первый международный кризис}~--- 1857 год.
		\item \textbf{Великая депрессия}~--- 1929--1933 гг., США.
		\item \textbf{Нефтяной кризис}~--- 1970-е годы.
		\item \textbf{Трансформационный кризис}~--- конец 1980-х -- начало 1990-х (постсоциалистические страны).
		\item \textbf{Мировой финансовый кризис}~--- 2008--2009 гг.
	\end{itemize}
	
	\subsection{Тенденции развития кризисов}
	\begin{enumerate}
		\item Кризисы приобрели мировой характер.
		\item Сокращается период между кризисами.
		\item Усиливается роль случайных факторов.
	\end{enumerate}
	
	\section{Факторы производства и факторные доходы}
	
	\subsection{Понятие факторов производства}
	
	\textbf{Факторы производства}~--- экономические ресурсы, используемые для создания товаров и услуг.
	
	\subsection{Основные факторы производства и доходы от них}
	
	\begin{enumerate}
		\item \textbf{Земля}~--- все природные ресурсы (земельные участки, леса, водные ресурсы, полезные ископаемые).
		
		\textit{Доход: \textbf{рента}.}
		
		\item \textbf{Труд}~--- физические и умственные способности людей, используемые в производстве.
		
		\textit{Доход: \textbf{заработная плата}.}
		
		\item \textbf{Капитал}~--- средства производства, созданные человеком:
		\begin{itemize}
			\item Физический (реальный) капитал~--- здания, оборудование, станки.
			\item Финансовый капитал~--- деньги, ценные бумаги.
		\end{itemize}
		
		\textit{Доход: \textbf{процент}.}
		
		\item \textbf{Предпринимательские способности}~--- способность организовать производство, принимать решения, рисковать.
		
		\textit{Доход: \textbf{прибыль}.}
		
		\item \textbf{Информация} (добавлен в конце XX века)~--- знания, данные, технологии.
		
		\textit{Особенность: неисчерпаемость.}
		
		\textit{Доход: \textbf{роялти}.}
	\end{enumerate}
	
	\subsection{Характеристики факторов производства}
	\begin{itemize}
		\item \textbf{Ограниченность}~--- все ресурсы имеются в ограниченном количестве.
		\item \textbf{Взаимозаменяемость}~--- возможность частичной замены одного фактора другим.
		\item \textbf{Взаимодополняемость}~--- необходимость сочетания факторов для производства.
	\end{itemize}
	
	\section{Функции рынка, его преимущества и недостатки, пути их преодоления}
	
	\subsection{Условия возникновения рынка}
	\begin{itemize}
		\item Общественное разделение труда.
		\item Экономическое обособление производителей.
		\item Развитие товарно-денежных отношений.
	\end{itemize}
	
	\subsection{Структура рынка}
	\begin{itemize}
		\item \textbf{Рынок товаров и услуг} (потребительский).
		\item \textbf{Рынок факторов производства} (включая рынок земли).
		\item \textbf{Финансовый рынок} (кредитный и фондовый).
		\item \textbf{Рынок труда}.
	\end{itemize}
	
	\subsection{Функции рынка}
	\begin{enumerate}
		\item \textbf{Посредническая}~--- связывает производителей и потребителей.
		\item \textbf{Ценообразующая}~--- устанавливает равновесную цену через механизм спроса и предложения.
		\item \textbf{Информационная}~--- предоставляет информацию о ценах, товарах, условиях продажи.
		\item \textbf{Регулирующая}~--- направляет ресурсы в наиболее востребованные отрасли.
		\item \textbf{Санирующая (оздоровительная)}~--- очищает экономику от неэффективных производителей.
		\item \textbf{Стимулирующая}~--- побуждает к внедрению новых технологий, снижению издержек.
	\end{enumerate}
	
	\subsection{Преимущества рынка}
	\begin{enumerate}
		\item Эффективное распределение ресурсов.
		\item Гибкость и быстрая адаптация к изменениям.
		\item Стимулирование научно-технического прогресса.
		\item Свобода выбора для производителей и потребителей.
		\item Разнообразие товаров и услуг.
	\end{enumerate}
	
	\subsection{Недостатки (провалы) рынка}
	\begin{enumerate}
		\item Не обеспечивает производство общественных благ (оборона, образование, здравоохранение).
		\item Не решает социальные проблемы (безработица, бедность, неравенство).
		\item Порождает социальное расслоение.
		\item Не учитывает внешние эффекты (загрязнение окружающей среды).
		\item Склонность к монополизации.
		\item Нестабильность (кризисы, инфляция).
	\end{enumerate}
	
	\subsection{Пути преодоления недостатков (роль государства)}
	
	\textbf{Механизмы государственного регулирования:}
	\begin{enumerate}
		\item \textbf{Правовое регулирование}~--- создание законодательной базы для функционирования рынка.
		\item \textbf{Фискальная политика}~--- использование налогов и государственных расходов для регулирования экономики.
		\item \textbf{Монетарная (денежно-кредитная) политика}~--- регулирование денежной массы и процентных ставок.
		\item \textbf{Антимонопольное регулирование}~--- предотвращение монополизации рынков.
		\item \textbf{Социальная политика}~--- защита малоимущих, пенсионное обеспечение, пособия.
		\item \textbf{Производство общественных благ}~--- финансирование обороны, образования, здравоохранения.
		\item \textbf{Экологическая политика}~--- установление экологических стандартов и штрафов.
		\item \textbf{Лицензирование}~--- контроль качества и безопасности продукции.
	\end{enumerate}
	
	\textbf{Вывод:} Современная экономика представляет собой смешанную систему, где рыночные механизмы дополняются государственным регулированием для преодоления провалов рынка и обеспечения социальной стабильности.
\end{document}
